\section{Where we stand: one network per interface is exact}
\label{sec:standing}

Part I trained on $85\%$ of an interface's (state, action) cells and scored the rest, because the question was
whether a network could generalise a rule it had seen only in part. That is not the question this programme
asks. Each interface gets its own network, trained on the whole of its spec; nothing is held out, because
nothing needs to be. The questions that remain are the ones a compiler is asked: is the output exact, is it
exact every time, and what does it cost. This section answers them for \nPlanned{} interfaces, \nNewDemos{} of them rebuilt
from the demo set of declarative dragging~\citep{dragology}, and then shows where the one-network-per-interface recipe stops.

\subsection{Trained on its whole spec, one network per interface is exact}
\label{sec:round2}

\paragraph{The interfaces.} \Cref{tab:round2} lists them. \nBefore{} were there before this round: Part I's
five hand-written specs (toggle, counter, both sliders, negotiate; its three screenshot-derived specs were never
committed and are not here) and two added since (\texttt{tracks}, a block on rails; \texttt{reorder},
Dragology's permutation demo). \nNewDemos{} are Dragology demos chosen to test what Part I's specs did not: rules whose guard reads other fields
(hanoi: only a top disk moves, never onto a smaller one); a two-dimensional reorder (kanban); a sliding puzzle
(fifteen); cars that block each other (rush hour); a rule over three objects at once (sokoban: the player, the
box ahead, the cell beyond); a list whose length changes (insert and remove); a reorder with a flag per item
(todo); meetings whose ends move without crossing (scheduler); a three-way \texttt{between} with a mode that
allows only the next place clockwise; and three continuous sliders that move together and stop together
(linked sliders). Hanoi also carries a \texttt{learn} hole: \texttt{hint}, the optimal next move toward the
goal peg, is given only by example, one per position, found by search when the spec is generated; the spec
writes the rules of the puzzle and shows, but does not write, how to solve it.

Each spec header records what was scaled down to fit enumeration (fifteen is $3\times2$, not $4\times4$;
sokoban is the top of Dragology's level; rush hour is $5\times5$ with six cars) and what Dragology does that
the spec language cannot yet say: drags that chain several snaps in one gesture, a snap radius, text, and
drawing anything but points. The eight specs written by coding agents have integration tests
(\texttt{crates/nerve/tests/demo\_*.rs}) showing that the drag futures the engine derives equal the list
Dragology's demo writes by hand (\texttt{moveItem(items, i, j)} for every $j$, the four directions a Sokoban
player may step, every stop a car can reach), from several starting states and, for fifteen, rush hour,
sokoban, insert/remove and the switch, from every valid one; hanoi, kanban, reorder and tracks have unit tests
of the same kind in \texttt{drag.rs}. The spec language has no list type, so a list is one position field per
item and an invariant that no two share a place; kanban's 360 boards live among $20{,}736$ slot tuples, and
\cref{sec:valid} is about that gap.

\begin{table}[t]
\centering
\caption{\textbf{Trained on its whole spec, one network per interface is exact.} Each row is one interface and
its own network, trained on every (state, action) cell with the configuration of Part I (300 epochs, AdamW,
weight decay $0.05$, factored head on) on a local RTX A4500; continuous specs train on $8{,}000$ sampled pairs
for $1{,}000$ epochs and are scored on $1{,}000$ fresh ones. \emph{States} counts states the invariants allow;
\emph{cells} counts legal (state, action) pairs; \emph{wall s} is wall-clock training time with several runs
sharing the GPU (\texttt{docs/paper2/timings.md}); \emph{exact} is the worst seed and \emph{violations} the
worst seed's rate of answers that break a rule. Rows are generated from \texttt{experiments/round2/results.jsonl}
and fill in as the sweep completes.}
\label{tab:round2}
\small
\setlength{\tabcolsep}{4pt}
\resizebox{\linewidth}{!}{%
\begin{tabular}{llp{4.6cm}rrrrrlr}
\toprule
interface & source & what it tests & states & cells & params & wall s & seeds & exact & viol. \\
\midrule
\input{table_round2_rows}% ends with \bottomrule
\end{tabular}}
\end{table}

\paragraph{Result.} Every finite interface is exact: $100\%$ of its cells, no violation, on every seed tried
(\nExactSpecs{} of \nFiniteSpecs{}; the seeds column says how many each has, three for all but the two
largest). The scheduler, the largest at \maxCells{} cells, was exact from its first evaluation. Hanoi reproduces all 81 solver examples on all three seeds while also being exact on its 729
written cells, so a hole and a rule can share one network of $105$k parameters without the examples being
drowned by the rules. The continuous interfaces are a different story. The two with one continuous field are
within $1\%$ of range on $99.6\%$ and $100\%$ of fresh transitions with $0.1\%$ and $0$ violations, at Part I's
tolerance under a longer recipe. The third, linked sliders, is the first interface in this programme that a
network trained on its whole spec does \emph{not} reproduce: $10.4\%$ of fresh transitions within $1\%$ of
range, a mean error of $7.5\%$ of range, and $89.6\%$ of answers breaking a rule. Its three continuous fields
move together and stop together at a track end, so the step is a descent over three coupled dimensions from
eleven seeds, and its \texttt{ensures} hold the group's shape to $0.001$, a thousandth of what the network's
tolerance allows, so an answer that is nearly right still breaks the rule. We leave the row in the table as it
is; \cref{sec:handoff} and \cref{sec:limits} draw the consequence, which is that continuous blocks must be
small and their hand-offs discrete. Nothing in \cref{tab:round2} needed a larger network than Part
I's, or a new idea; it needed exhaustive coverage of the cells and the discipline below. The whole sweep is
\totalGpuMinutes{} wall-clock minutes with several runs sharing the GPU, and the largest interface
(\maxCellsSpec, \maxCells{} cells) took over an hour of it.

\paragraph{Discipline.} Part I's own review found specs, configurations and logs behind its tables
uncommitted, and this round found one more failure of the same kind: served models whose records had been
patched by hand after training. Every model in \cref{tab:round2} is now produced by one command
(\texttt{experiments/round2/run.sh}, committed with this paper), writes its own \texttt{model.json} once and
is never edited, is scored by a second command into an append-only log, and the set the lab serves is chosen
from that log by a third. A seed is a folder; a rerun is a new row that replaces the old one in the summary
rather than a second vote. We mention this because the next sections propose machines of many networks, and a provenance discipline
that is merely good practice for one network is a precondition for debugging twenty.

\subsection{Answering only from states the spec allows removes rule breaks without retraining}
\label{sec:valid}

Reorder was the first interface whose network broke rules when data was scarce. Trained on $30\%$ of its
cells it answered $61.7\%$ of the rest correctly and broke an invariant on $26.6\%$ of them: two items in one
slot. The trainer scored every one of the $3{,}125$ position tuples as a candidate successor, and only $120$
of them are orders. The spec knows which $120$; the network was never told.

The fix is to restrict the argmin to the states the invariants allow and give every other state infinite
energy. This is the candidate-set rule of our declarative-dragging note (a state outside the derived futures
costs infinite energy), in the spirit of Part I's Theorem~1, applied to the learned term's own argmin; and it
is the first wiring primitive of \cref{sec:blocks}: a block's output is typed by its invariant. On the same
trained networks, with no retraining (every number here is from \texttt{experiments/round1b/}), reorder went from $61.7\%$
to $95.3\%$ at $70\%$ held out and from $29.8\%$ to $68.7\%$ at $90\%$, with violations falling from $27\%$ and
$63\%$ to $0.1\%$ and $5\%$ (the residue is an \texttt{ensures}, not an invariant: a legal order in which the
moved item is not where it was dropped). Retraining over the $120$ orders added nothing ($95.2\%$ against
$95.3\%$): the network already ranked real orders correctly, and impossible tuples were beating them at the
argmin. Interfaces with no invalid states score identically either way, by construction. Valid candidates are
the default from here on; \texttt{nerve eval --all-states} scores the network on its own when a paper needs
that number, and we report both wherever they differ.

\subsection{Drags are searches over futures the spec derives}
\label{sec:drags}

Every interface in \cref{tab:round2} is operated by dragging in the lab that serves it. A drag is one
minimisation per pointer frame over the futures the spec allows from the current state: the pixel gap from
the pointer to where the grabbed element would land, plus the network's energy as a tie-breaker (a candidate's
energy less the minimum over that action's candidates, infinite outside the set), with
Dragology's operators (\texttt{closest} with stickiness, \texttt{between}, \texttt{vary} through the spec's
own clamps, \texttt{far} with hysteresis) and two added for the rails demo: \texttt{each}, a \texttt{closest}
whose every branch is a \texttt{vary}, and \texttt{floating}. The futures are never listed; they are derived by
asking the oracle which actions reach which values, so a drag cannot offer an illegal drop (a buried Hanoi disk
has nowhere to go; a car stops at the first car in its way). Two facts matter for what follows. The network's
term is exactly zero on every finite interface, because the network and the spec agree on every cell, and small
on the continuous ones; in a drag the network's job is ties and holes, not geometry. And a drag is already a
search that re-plans under the hand at every frame, which \cref{sec:runtime} will lean on when the thing being
dragged is part of a running machine.

\subsection{One network per interface hits an enumeration wall}
\label{sec:wall}

\begin{figure}[t]
\centering
\input{fig_cost}
\caption{\textbf{One network per interface hits an enumeration wall.} Wall-clock seconds for 300 epochs
against the number of (state, action) cells, finite interfaces of \cref{tab:round2}, with several runs sharing
the GPU. Cost grows roughly with the number of cells; the full fifteen puzzle has about $10^{13}$ states and
no cell count that fits on this axis.}
\label{fig:cost}
\end{figure}

The trainer enumerates the state space and trains on every cell, and its wall-clock cost grows roughly with
the number of cells (\cref{fig:cost}; the candidates' embeddings are encoded once per step, so the state count
enters weakly): hanoi's 729 cells train in six minutes, kanban's $17{,}280$ in over an hour, with several runs
sharing the GPU. The interfaces we scaled down to fit are the ones
Dragology is proud of at full size: the $4\times4$ fifteen puzzle has $16!/2\approx10^{13}$ states, a real
Sokoban level has millions, and anything with a list of lists has no fixed size at all. Two exits exist. One is
engineering: score only a gesture's futures rather than every state, which turns the wall into a per-action
cost and is on the plan (\cref{sec:futures}). The other is the subject of this paper: stop building one network
per machine. A million-state odometer is six ten-state digits and five carry wires, and its network should be
one ten-state network used six times.
